
\section{Design Rationale}
\label{sec:design_rationale}

The design of the prototype is guided by modularity, configurability, and observability. Since the objective of the project is to support controlled ZK-rollup benchmarking rather than production deployment, the architecture separates responsibilities across components and uses established software design patterns to make individual stages easier to modify, test, and measure.

\begin{itemize}
    \item \textbf{Strategy Pattern:} 
    The Sequencer uses the Strategy pattern to support multiple transaction ordering policies, such as FCFS, FeePriority, and TimeBoost, without changing the core batching logic. This allows scheduling policies to be swapped during experiments while keeping the rest of the pipeline unchanged.

    \item \textbf{Pipeline Pattern:} 
    The Executor and Submitter follow a pipeline-style structure where batches pass through clear processing stages. In the Executor, this includes validation, state-transition execution, trace generation, and prover invocation. In the Submitter, this includes payload preparation, data availability formatting, gas estimation, and Layer-1 submission.

    \item \textbf{Saga Workflow Pattern:} 
    The Submitter uses a Saga-style workflow to manage reliable Layer-1 submission. This is necessary because interactions with the Layer-1 RPC can fail due to nonce conflicts, temporary node instability, or transaction submission errors. The Saga workflow supports retry handling, idempotency, and recovery without requiring the Sequencer or Executor to stop processing new batches.

    \item \textbf{Observer Pattern:} 
    The Sequencer uses an event-listening mechanism to observe relevant Layer-1 bridge events, such as forced transactions or deposit-related updates. This separates event monitoring from the main transaction batching logic and keeps the Sequencer extensible for future rollup features.
\end{itemize}

These design choices support the central research goal of the project. By isolating scheduling, execution, proving, submission, and monitoring responsibilities, the prototype allows individual configuration variables to be changed and evaluated independently. This makes the framework suitable for studying throughput, latency, finalization time, data availability cost, and prover-related bottlenecks under controlled experimental conditions.

